\subsection{p-环的存在性与唯一性}

命题 4．——设 C 和 C' 是两个 p-环，满足 \(l(\mathbf{C}) \geqslant l(\mathbf{C}')\)，并设 \(\pi\)（分别为 \(\pi'\)）是从 C（分别从 C'）到 \(\kappa_{\mathbf{C}}\)（分别到 \(\kappa_{\mathbf{C}'}\)）的典范同态。设 \((x_{\lambda})_{\lambda \in \Lambda}\)（分别为 \((x_{\lambda}')_{\lambda \in \Lambda}\)）是 C（分别 C'）中的元素族，其在 \(\pi\)（分别 \(\pi'\)）下的像构成 \(\kappa_{\mathbf{C}}\)（分别 \(\kappa_{\mathbf{C}'}\)）的 p-基。设 v 是从 \(\kappa_{\mathbf{C}}\) 到 \(\kappa_{\mathbf{C}'}\) 的同构，满足 \(v(\pi(x_{\lambda})) = \pi'(x_{\lambda}')\)，对每个 \(\lambda \in \Lambda\) 均如此。于是存在唯一一个从 C 到 C' 的同态 u，使得 \(v \circ \pi = \pi' \circ u\)，且 \(u(x_{\lambda}) = x_{\lambda}'\)，对每个 \(\lambda \in \Lambda\) 均如此。它是满射。若 \(l(\mathbf{C}) = l(\mathbf{C}')\)，它是同构。

证明 u 的存在性。令 A 为 C × C' 的子环，由对 \((x, x')\) 构成，这些对满足 \(v(\pi(x)) = \pi'(x')\)。映射 \((x, x') \mapsto \pi(x)\) 是从 A 到 \(\kappa_{\mathbb{C}}\) 的满射环同态。它的核 m 等于 \(p\mathbb{C} \times p\mathbb{C}'\)，因而是 A 的极大理想。\(\mathbb{C} \times \mathbb{C}'\) 的拓扑子空间 A 在 \(\mathbb{C} \times \mathbb{C}'\) 中是闭的，因而完备；由 \(\mathbb{C} \times \mathbb{C}'\) 在 A 上诱导的拓扑是 m-进拓扑。因此 A 是以 m 为极大理想的分离且完备的局部环（III，§ 2，第 13 号，命题 19）。对于每个 \(\lambda \in \Lambda\)，根据假设有 \((x_{\lambda}, x_{\lambda}') \in A\)；若 \(\xi_{\lambda}\) 是 \((x_{\lambda}, x_{\lambda}')\) 模 m 的类，则族 \((\xi_{\lambda})_{\lambda \in \Lambda}\) 是域 A/m 的一个 \(p\)-基。根据第 2 号的定理 1，存在唯一一个 A 的科恩子环 \(\mathbb{C}''\)，包含 \((x_{\lambda}, x_{\lambda}')\)（对每个 \(\lambda \in \Lambda\)）。有 \(l(\mathbb{C}'') = l(\mathbb{C}) \geqslant l(\mathbb{C}')\)。在 \(\mathbb{C}''\) 上将 \(\mathbb{C} \times \mathbb{C}'\) 到 C 的投影限制后得到同态 \(h: \mathbb{C}'' \to \mathbb{C}\)，它诱导出从 \(\kappa_{\mathbb{C}''}\) 到 \(\kappa_{\mathbb{C}}\) 的同构。根据第 2 号的命题 2 c)，h 是从 \(\mathbb{C}''\) 到 C 的同构。同理，限制 \(h'\) 是在 \(\mathbb{C}''\) 上由 \(\mathbb{C} \times \mathbb{C}'\) 到 C' 的投影得到的满射同态，其从 \(\mathbb{C}''\) 到 C'。因此，\(\mathbb{C}''\) 是从 C 到 C' 的满射同态 \(u = h' \circ h^{-1}\) 的图像，并且显然有 \(v \circ \pi = \pi' \circ u\)，\(u(x_{\lambda}) = x_{\lambda}'\)，对每个 \(\lambda \in \Lambda\) 均如此。此外，若 \(l(\mathbb{C}) = l(\mathbb{C}')\)，u 是同构。

证明 \(u\) 的唯一性。设 \(u_{1}\) 是从 C 到 C' 的同态，满足 \(v \circ \pi = \pi' \circ u_{1}\) 且 \(u_{1}(x_{\lambda}) = x_{\lambda}'\)，对每个 \(\lambda \in \Lambda\) 均如此，并令 \(\mathbf{C}_{1}\) 为 \(u_{1}\) 的图。显然，\(\mathbf{C}_{1}\) 是 A 的科恩子环，包含 \((x_{\lambda}, x_{\lambda}')\)，对每个 \(\lambda \in \Lambda\) 均如此，从而 \(\mathbf{C}_{1} = \mathbf{C}''\)（第 2 号的定理 1），最后 \(u_{1} = u\)。

命题 5．——设 k 是特征为 p 的域，且 n 是满足 ≥ 1 的整数或 + ∞。存在一个长度为 n 且剩余域同构于 k 的 p-环。

Witt 向量环 \(\mathbf{W}(k)\) 的系数在 \(k\) 中，是分离且完备的整局部环，其剩余域同构于 \(k\)（\(\S\) 1，第 8 号，命题 8），且有 \(p\)．\(1_{\mathbf{W}(k)} \neq 0\)（同上，公式（52））。令 C 为 \(\mathbf{W}(k)\) 的一个科恩子环（第 2 号，定理 1）。则 C 是一个 \(p\)-环，长度为 +\(\infty\)，其剩余域同构于 \(k\)；若 \(n\) 是满足 \(\geqslant 1\) 的整数，则商 \(\mathrm{C}/p^{n}\mathrm{C}\) 是一个 \(p\)-环，长度为 \(n\)，其剩余域同构于 \(k\)。

注记．——1) 令 \(n\) 为满足 \(\geqslant 1\) 的整数，并令 S 为一个 \(p\)-基（属于 \(k\)）。可以证明，\(\mathbf{W}_{n}(k)\) 中由 \(\mathbf{W}_{n}(k^{p^{n}})\) 及元素 \([\xi, 0, ..., 0]\)（\(\xi \in S\)）生成的子环，是一个 \(p\)-环，长度为 \(n\)，其剩余域同构于 \(k\)（参见 p. 72，习题 10）。

\begin{enumerate}
\def\labelenumi{\arabic{enumi})}
\setcounter{enumi}{1}
\tightlist
\item
  读者将在附录中找到命题 5 的一个证明；该证明既不使用 § 1 的结果，也不使用科恩子环的存在定理（第 2 号，定理 1）。
\end{enumerate}

推论．——设 C 是长度有限的 p-环，长度为 n。存在一个无限长度的 p-环 C'，使得 C 同构于 \(C'/p^{n}C'\)。

根据命题 5，存在一个无限长度的 \(p\)-环 \(C'\)，使得 \(\kappa_{C'}\) 同构于 \(\kappa_{C}\)。于是 \(C'/p^{n}C'=C_{n}'\) 是一个 \(p\)-环，长度为 \(n\)，且域 \(\kappa_{C_{n}'}\) 同构于 \(\kappa_{C'}\)，因而同构于 \(\kappa_{C}\)。根据命题 4，环 \(C\) 与 \(C_{n}'\) 因而同构。

